\section{The stability dichotomy and the final algorithm}
\label{asm:section}

We now combine the two constructions. The distinction between the
cases is used only to analyze the candidates: the algorithm computes
all candidates without knowing any optimum cost.

\subsection{Choosing the constants}

Take $C_0$ from Equation~\eqref{loc:main-count-constant}, put
$T=200C_0$, and let $c=c(T)>0$ be given by
Equation~\eqref{str:constants}. Fix the rational constants
\begin{equation}\label{asm:constants}
 \gamma=\min\{3/100,c/100\},\qquad
 \xi=\varepsilon_0=\gamma/16,\qquad
 \sigma=\gamma/4.
\end{equation}
In particular $0<\gamma<1$. The choice of $\varepsilon_0$ fixes the
polynomial distance bound in Lemma~\ref{pre:rounding} before the
external surplus construction is invoked. Choose $v$ as in
Equation~\eqref{sur:parameter-choices}, then choose $N_*$ and the
common integer $L$ as in Equation~\eqref{sur:size-threshold} and the
paragraph following it. Thus
\[
 L=O(\log(N+2)),\qquad L\ge1,
\]
where all hidden constants are fixed independently of the input.
There is no dependence of $C_0,T,c,\gamma,\xi$, or $\varepsilon_0$
on $L$. This order also accommodates the fixed coefficient in the
rounded distance bound.

\begin{lemma}[A sensitive cardinality]\label{asm:sensitive}
Suppose $k>2L$, $0<\xi\le1$, and
\[
 \operatorname{opt}_{k-2L}>(1+\xi)\operatorname{opt}_k.
\]
Put $\beta=\xi/(4L)$. There is an integer $h'$ with
$k-2L<h'\le k$ such that
\begin{equation}\label{asm:sensitive-ratios}
 \operatorname{opt}_{h'-1}\ge(1+\beta)\operatorname{opt}_{h'},
 \qquad
 \operatorname{opt}_{h'}\le(1+\xi)\operatorname{opt}_k.
\end{equation}
For an optimum $O$ of size $h'$, of cost
$P=\operatorname{opt}_{h'}$, every comparison cluster satisfies
\begin{equation}\label{asm:separation}
 n_i\operatorname{sep}_i\ge\beta P,
 \qquad \operatorname{sep}_i=d(i,O\setminus\{i\}).
\end{equation}
\end{lemma}

\begin{proof}
All optimum costs in the integral model are positive. Since
$\log(1+\xi)\ge\xi/2$ for $0\le\xi\le1$ and
$\log(1+\beta)<\beta$,
\begin{equation}\label{asm:product-bound}
 (1+\beta)^{2L}<\exp(\xi/2)\le1+\xi.
\end{equation}
The product of the $2L$ ratios
$\operatorname{opt}_{r-1}/\operatorname{opt}_r$, for
$k-2L<r\le k$, is greater than $1+\xi$. Hence at least one
ratio is at least $1+\beta$. Choose the largest such index $h'$.
Every later ratio is less than $1+\beta$, and their product is
at most $(1+\beta)^{2L}<1+\xi$ (the empty product is one).
This proves both inequalities in
Equation~\eqref{asm:sensitive-ratios}.

We have $h'\ge k-2L+1\ge2$. For $i\in O$, reroute all its
clients to a nearest other member of $O$ and keep the remaining
assignments. This constructs an $(h'-1)$-center solution of cost
at most $P+n_i\operatorname{sep}_i$. The first inequality in
Equation~\eqref{asm:sensitive-ratios} gives
$P+n_i\operatorname{sep}_i\ge(1+\beta)P$, proving
Equation~\eqref{asm:separation}. In particular no cluster of this
optimum is empty.
\end{proof}

\subsection{Candidates on the bounded integral instance}

\begin{proposition}[Integral-instance algorithm]\label{asm:integral}
With the constants above, for every $0<\zeta<1$ there is a randomized
algorithm which always returns at most $k$ facilities and cost at
most $5\operatorname{opt}_k$, and with probability at least
$1-\zeta$ returns cost at most
\begin{equation}\label{asm:integral-guarantee}
 (2-\gamma)(1+\xi)\operatorname{opt}_k.
\end{equation}
Its running time is polynomial in $N$ and $\log(2/\zeta)$ for the
fixed constants and the fixed polynomial distance bound.
\end{proposition}

\begin{proof}
Compute a width-$5$ local optimum $S_k$ of size $k$ and retain it
as a candidate. Lemmas~\ref{loc:descent} and~\ref{loc:swap} give
polynomial time and $\operatorname{cost}(S_k)\le5\operatorname{opt}_k$.
Every other retained candidate will also satisfy the budget, and we
return the least-cost candidate.

If $k\le2L$, apply Theorem~\ref{comp:recovery} to $S_k$ with
the supplied count bound $M_0=2L$ and failure parameter $\zeta$.
For its analysis, take an optimum of size $k$ and declare every
cluster bad. Then $m=k\le M_0$, and the proxy condition is vacuous.
Recovery gives cost at most $1.97\operatorname{opt}_k$ with
probability at least $1-\zeta$. Since
$1.97\le2-\gamma$, this implies
Equation~\eqref{asm:integral-guarantee}, including the case $k=1$.

Suppose now that $k>2L$. Compute the candidate of
Proposition~\ref{sur:low-sensitivity} with target $h=k-L$.
In addition, for every integer $h'$ with $k-2L<h'\le k$, compute
a width-$5$ local optimum $S_{h'}$ of size $h'$ and apply
Theorem~\ref{comp:recovery} with
\begin{equation}\label{asm:count-bound}
 M_0=\left\lceil\frac{4C_0L}{\xi}\right\rceil,
 \qquad\text{failure parameter }\zeta.
\end{equation}
These calls require only the computed local sets, $M_0$, and $\zeta$.
They enumerate their unknown comparison costs internally.
All their feasible outputs have at most $h'\le k$ centers.

If $\operatorname{opt}_{k-2L}\le(1+\xi)\operatorname{opt}_k$,
Proposition~\ref{sur:low-sensitivity} gives cost at most
\[
 (2-c/100)(1+\xi)\operatorname{opt}_k
 \le(2-\gamma)(1+\xi)\operatorname{opt}_k.
\]
This conclusion is deterministic.

Otherwise, let $h'$ be supplied by Lemma~\ref{asm:sensitive},
take an optimum $O$ of size $h'$, and put
$P=\operatorname{opt}_{h'}$, $b=\beta P$.
The local set $S_{h'}$ has no surplus relative to $O$.
Equation~\eqref{asm:separation} says that no center is short
at scale $b$, so the left side of the scale condition in
Theorem~\ref{loc:anchoring} is zero. With tolerance $1/100$,
that theorem supplies distinct good proxies with total assignment
cost at most $P_g+P/100$, and
\[
 m\le C_0P/b=C_0/\beta=4C_0L/\xi\le M_0.
\]
The recovery call at this particular $h'$ therefore has, with
probability at least $1-\zeta$, cost at most
\[
 1.97P\le1.97(1+\xi)\operatorname{opt}_k
 \le(2-\gamma)(1+\xi)\operatorname{opt}_k.
\]
The algorithm need not recognize this index or its success event:
its returned minimum is no worse than this candidate. Only this
one call must succeed, so there is no union-bound loss over the
other indices.

Both supplied count bounds are $O(\log(N+2))$. There are at most
$N$ cardinality choices; the local searches and the surplus
construction are polynomial, and the recovery calls have polynomial
time by Theorem~\ref{comp:recovery}. Keeping $S_k$ proves the
unconditional factor-five bound.
\end{proof}

\subsection{General metrics and the expected guarantee}

\begin{proof}[Proof of Theorem~\ref{main:theorem}]
Handle the elementary cases and zero optimum as in
Section~\ref{pre:section}. Apply Lemma~\ref{pre:rounding} with
the fixed $\varepsilon_0=\gamma/16$ and the algorithm of
Proposition~\ref{asm:integral}. The successful approximation factor
on the original rational metric is at most
\[
 (2-\gamma)(1+\xi)(1+\varepsilon_0)
  =(2-\gamma)(1+\gamma/16)^2
  \le2-\gamma/2.
\]
For the last inequality, expand the product:
\[
 (2-\gamma)(1+\gamma/16)^2
 =2-\frac34\gamma-\frac{15}{128}\gamma^2
                  -\frac1{256}\gamma^3
 \le2-\gamma/2.
\]
The transferred unconditional bound is
$5(1+\varepsilon_0)\operatorname{opt}_k<6\operatorname{opt}_k$.

For the prescribed fixed $a>0$, choose a rational failure parameter
\[
 0<\zeta\le\min\{(N+2)^{-\lceil a\rceil},\gamma/24\}.
\]
It has polynomial bit length and
$\log(2/\zeta)=O_a(\log(N+2))$ with fixed constants.
The least original-cost candidate over all rounding guesses is
no worse than the candidate for the correct guess. Hence its
probability of exceeding $(2-\gamma/2)\operatorname{opt}_k$
is at most $\zeta\le(N+2)^{-a}$.
Since $\sigma=\gamma/4$, this proves the asserted probability bound.
Finally, split expectation according to that successful event and
use the unconditional cost bound:
\[
 \mathbb E\operatorname{cost}(S)
 \le (2-\gamma/2)\operatorname{opt}_k
              +6\zeta\operatorname{opt}_k
 \le(2-\gamma/4)\operatorname{opt}_k
 =(2-\sigma)\operatorname{opt}_k.
\]
Every candidate is a real facility set of size at most $k$.
All enumerations and arithmetic have polynomial bit complexity,
as established in the reduction, sampling, and completion proofs.
This completes the theorem.
\end{proof}
